\chapter{Theoretical Framework}
\thispagestyle{empty}
This chapter combines significant findings from various literature and studies from tourism, exploration of the traveling salesman problem (TSP) in itinerary planning, and methods for recommendation for itinerary generation. This comprehensive review serves as the basis for the development of the Lakad application by identifying relevant factors such as optimization algorithms for TSP in itinerary planning and factors to consider when creating personalized recommendation features of the system.

\section{Relevant Theories}
Key features of Lakad rely on established theories of mathematics. This section examines the foundational theories and explains how they are used and related to the concept of itinerary optimization and personalized generation.

\subsection{Graph Theory}
Graph Theory is a branch of mathematics that is still relatively young but is developing quickly. It originated in the 18th century when Leonhard Euler solved the Königsberg Bridge Problem, which established the first challenge of that academic field. Euler introduced the method of using vertices and edges to show the real-world connections, establishing the foundation of all network studies and relations. In the early 20th century, mathematicians like Kuratowski, Wagner, and Whitney expanded the field through studying planar graphs that led to the development of basic concepts that became Graph Minor Theory, which established connections between graph structures and both geometry and topology. The proof of 1976 Four color theorem increased the popularity of Graph Theory and demonstrated its ability to solve difficult mathematical challenges \parencite{Carmesin2022}.

A graph is a set of vertices and edges, where each edge connects exactly two vertices. Topologically, a graph can be thought of as a one-dimensional simplicial complex \parencite{Carmesin2022}. A graph $G=(V,E)$ is composed of a set of vertices $V$ and a set of edges $E$. There are two types of graphs, directed graphs, which have edges that can only go in one direction, and undirected graphs, where edges allow movement in both directions. Additionally, they can be weighted, it means that each edge has a number that represents different factors like cost, time, and distance.

In this study, the concept of graph theory was applied to represent the points of interest (POI) as nodes and the distance between every connected vertex as the edges. This graphical representation of cities and their distances is one of the common applications of graph theory and allows access to powerful algorithms and methods for optimization.

\subsection{Traveling Salesman Problem}

This is a specific problem under the field of graph theory. According to \textcite{Pop2024}, the Traveling Salesman Problem (TSP) has been in the history of combinatorial optimization since 1930, and was properly provided a mathematical formulation by Merrill M. Flood. TSP became very popular since it is a widely investigated optimization problem, often serving as a benchmark for modern optimization algorithms. The problem states that, given a set of cities, the salesman’s goal is to find the shortest possible route that visits each city exactly once.  According to \textcite{Wu2020}, TSP is considered a non-deterministic polynomial time problem or simply an NP-hard problem, meaning it is hard to solve efficiently and finding the exact solution takes a lot of time. Hence, one of the best way to solve it is through heuristic algorithm, where the best solution is not necessarily the optimal but the solution that takes reasonable of amount of time to solve \parencite{Sabery2023}.

A Study by \textcite{Violina2021}, analyzes two exact algorithms to solve TSP, the brute force algorithm and the branch and bound algorithm. They conclude that solving TSP with exact algorithm results in optimal outcome but less effective. On the other hand, independent studies by \textcite{Wu2020} and \textcite{Wadi2025}, utilizes metaheuristic to solve TSP, and found that it is much better than using exact algorithms. It can easily be seen how the idea of TSP can be applied to itineraries. Each location of itineraries are the nodes and the distances between them are the edges. The goal of the system for the itinerary is to find the route which covers the least distance \parencite{Pop2024}. These concepts of TSP are used in the route optimization feature of the system, where the system will find the optimal route for the generated itinerary.

\subsection{Orienteering Problem}
The Orienteering Problem (OP) is a routing problem that aims to maximize the total rewards collected along a route within a given travel budget \parencite{yu2022}. Similar to this, \textcite{morandi2024} explains OP as the task of planning a route for one vehicle that has to follow a travel budget. The problem is depicted as a graph where the roads have distances and the locations have points or reward values. The goal in this problem is to create a path that starts and ends at the depot, stays within the travel limit, and collects as many rewards as possible from the places visited. Similarly, it originated from a sport of the same name, where participants visit check-points with pre-determined scores, in an attempt to maximize their total score within a specific time \parencite{Lim2019}. \textcite{pvenivcka2019} described OP as a type of problem in operations research, introduced to it in 1984 by Tsiligirides. They then explained that OP is a problem that combines the combinatorial optimization problem Knapsack Problem (KP) with TSP. These two work separately as TSP, given a set of customers, seeks to find the sequence in which to visit selected customers, constrained within a budget, while also minimizing the tour path. The subset of selected customers is selected by the KP, where it maximizes the collected profit based on the selection of customers to be visited within the budget constraints. This was used in the itinerary generation feature of the system where the user provides input for the time or distance constraint to generate an itinerary for the user. This is supported by two independent studies by \textcite{Yochum2020} that utilize Adaptive Genetic Algorithm with Dynamic Mutation and Crossover Probabilities and, \textcite{Tenemaza2020} by modeling the Tourist Trip Design Problem (TTDP) as a Time-Dependent Orienteering Problem with Time Windows (TDOPTW) in their itinerary planning systems.

\subsection{Simple Additive Weighting}

Simple Additive Weighting (SAW), is one of the most widely used methods for Multi-Criteria Decision Making (MCDM). It is used to evaluate and rank a set of alternatives based on multiple and often conflicting criteria. SAW works by calculating a single performance score for each alternative by summing the weighted performance ratings of those alternatives across all criteria \parencite{Hamed2023}.

The SAW method uses a defined procedure to transform various assessment criteria into a standard measurement system. The data requires normalization as the initial step because decision matrices contain different measurement systems which require aggregation of multiple criteria \parencite{Hamed2023}. \textcite{Malefaki2025} recommends the use of min-max normalization for SAW because this method maintains data integrity under minor data changes while detecting major trend changes. The SAW process requires the establishment of criterion weights which determine their decision-making significance. The total weight of these elements must equal 1 (or 100 percent). Decision makers can choose between two methods for weight assignment, which include objective weighting based on data analysis or subjective selection based on their personal preferences and expertise \parencite{Odu2019}. The calculation of the final score for each alternative involves multiplying the normalized ratings by their weights and then adding the results together. The optimal solution definition depends on which alternative achieves the highest total score or the lowest total score according to the specified objective \parencite{Hamed2023}.

SAW was used in this study to choose the most effective route optimization algorithm. The selection criteria assess both accuracy and speed. The algorithms underwent assessment using these two criteria before the selection process determined which algorithm would achieve the best score for implementation in the system.

\subsection{Technology Acceptance Model}
The Technology Acceptance Model (TAM) is one of the well-known theories in explaining and predicting users of information systems. The Theory was developed by Fred D. Davis in 1989 to explain the psychological factors that determine whether the individual will accept and use the new technology. \textcite{Marikyan2025TAM} state that the model derives from the Theory of Reasoned Action (TRA) which demonstrates how the person’s attitude and their perception of social norms affect their intentions to behave in particular ways. TAM changes the original model by concentrating on the variables that determine the users’ attitude towards a specific system in their technological environment.

The model provides a theoretical framework for assessing how likely the technology will be adopted in different situations. The model has become a widely used tool because it offers simple yet powerful explanatory capabilities which identify major psychological factors that affect users’ adoption behavior. Through this framework, researchers and developers can analyze how users perceive a system and identify what or where the modification can be made to enhance acceptance and satisfaction \parencite{Marikyan2025TAM}. TAM assumed that the acceptance of technology is primarily determined by an individual’s behavioral intention (BI), which is shaped by two cognitive responses: Perceived Usefulness (PU) and Perceived Ease of Use (PEOU) \parencite{Marikyan2025TAM}. According to \textcite{Aburbeian2022}, these two constructs are considered the most critical variables that can influence the use or rejection of the new technology.


\subsection{Software Product Assessment Framework (ISO/IEC 25010:2023)}

ISO/IEC 25010:2023 is part of the Systems and Software Quality Requirements and Evaluation (SQuaRE) series, and it identifies the models to be used for describing and evaluating the quality of software and system products. It provides common terminology and concepts that enable developers, consumers, and evaluators to specify quality requirements and to determine whether a product meets the requirements. The standard emphasizes that both functional and non-functional aspects of performance should be reflected in the defined characteristics and sub-characteristics used in the evaluation of the quality of the product \parencite{ISO25010Web,ISOIEC25010_2023}.

The model states that the product quality model consists of nine primary characteristics: functional suitability, performance efficiency, compatibility, interaction capability, reliability, security, maintainability, flexibility, and safety. Each characteristic is made up of sub-characteristics which describe quantifiable components of quality, like operability, scalability, coexistence, correctness, and completeness. The 2023 revisions include revised terminologies in addition to the replacement of construct names such as usability with interaction capability, and portability with flexibility. The new system ensures that modern software products are evaluated consistently to meet user expectations for performance, reliability, and safety throughout their life cycle, under the new theoretical basis in order to keep up with changing technological context.


\section{Review of Related Literature}
In this section, relevant studies regarding the developed application are presented. Covering topics related to tourism in Bulacan, techniques for optimization of itineraries and itinerary generation, along with several related approaches in personalized travel planning. The review of these studies established the foundation of the system, and the selection for the most suitable methods for mobile itinerary generation is justified.

\subsection{State of Tourism in Bulacan}

According to \textcite{Canet2024P}, Bulacan’s tourism development is highly influenced mainly by accessibility, infrastructure, along with the promotion and preservation of cultural and natural heritage. Although the rich history and heritage of Bulacan holds a great potential for its tourist economy, its growth is often limited by issues such as lack of promotion. Canet and Sunpongco (2025) looked for the reasons as to why some of Bulacan’s historical and cultural heritage destinations remain neglected in spite of the province's efforts. Their study revealed that the already-popular landmarks, mostly in the capital city of Malolos, attracts most of the would-be-visitors while attractions from other municipalities and cities in Bulacan yield relatively lower visits, reflecting the uneven promotion of Bulacan’s tourist sites.

This imbalance aligns with the findings of \textcite{Canet2024P}, implying that the issue regarding weak promotion is true. To help resolve this issue \textcite{Canet2025S} also examined promotional strategies currently implemented by tourism offices, where they found that online promotions such as social media campaigns, and partnering with bloggers and vloggers prove to be the most effective methods when engaging with audiences of the younger generation. Although brochures and community-driven activities still does play important roles, digital strategies are better suited for the \textit{Gen Z} and Millennial demographics. \textcite{Canet2025S} then emphasized the importance of collaborating with local government units, and other tourism-related communities to ensure consistent promotion and the sustenance of better tourism development.

The local government of Bulacan recognized the importance of tourism within their province, as proven by how historical and culturally rich it is. With this in mind, the Provincial Government of Bulacan, in 2023, launched the Bulacan Pamana Pass. This tourism initiative focuses on Bulacan’s local culture and the conservation of its heritage along with several community engagement practices. Bulacan Pamana Pass is a handed-out booklet consisting of twenty historical and cultural landmarks located within the province. Several of these landmarks hold significant contributions to Philippines’ history, such as the Barasoain Church and Casa Real de Malolos. It also includes religious and non-religious structures as landmarks and tourist attractions. The pass serves not only as a promotional strategy, but also offers a structured approach to travelling or visiting the province \parencite{Velasco2023}. Bulacan’s Pamana Pass strengthens the preservation of heritage, and promotional efforts of the province when it comes to tourism. However, it does not provide the optimized sequencing of travel destinations for visitors, a task that is still left to the tourists to plan on their own. This gap is what needs a more robust solution, such as the need for systems that not only can identify attractions but also generate efficient itineraries tailored to tourists’ preferences.

\textbf{Digital Adoption of Tourism in Bulacan.} Digital strategies also play a critical role in Bulacan’s tourism promotion, \textcite{Canet2023SanRafael}, studied how social media content is disseminated to the users in San Rafael, Bulacan and found that 57 percent agreed that social media helped them determine their next choice of destination and 61 percent agreed that social media is a good strategy to make more tourist locations known to people. In line with this, \textcite{DelaCruz2022} also surveyed 100 tourists that visited Dona Remedios Trinidad, Bulacan to determine the effectiveness of using social media in promoting a tourist destination and asked questions ranging from how social media influences traveler’s decisions to the problems faced by tourists when using social media as a guide for tourism. They found that most travelers agreed that social media does influence the choice of destination of travelers. While problems reported by tourists is the spread of false or misleading information, different expectations, and inaccurate photos / videos to different tourist destinations. Clearly, the importance and role of social media in making tourist locations well-known is irreplaceable, aligning with the findings of \textcite{Canet2025S}. However, it is also prone to misleading information, highlighting the need for verified and curated information of tourist attractions in Bulacan.

\subsection{Metaheuristic Algorithm for Route Optimization}

Route optimization can be modeled as a TSP, where it seeks the shortest possible route that visits a set of cities (nodes) exactly once. There are many algorithms that provide the solution to the TSP, one of those algorithms is the exact method. According to \textcite{Violina2021}, exact methods such as Brute Force (BF) and Branch and Bound (BB) algorithms are used to solve TSP. These methods involve evaluating each possible solution one at a time until all of them are explored, then comparing each solution and then selecting the smallest one. The exact algorithm guarantees the best solution, however it becomes inefficient when the number of cities becomes large.

Another approach in solving TSP is through metaheuristic algorithms. These are often referred to as a Nature-inspired Optimization Algorithm, it is designed to find the optimal or near-optimal solutions to complex problems by mimicking the behavior of natural systems, processes, or phenomena \parencite{Sabery2023}. Generally, metaheuristic algorithms begin with an initial solution and follow a set of heuristic principles to explore the search space and iteratively improve it. These guidelines help the algorithm avoid local optima and locate globally optimal or nearly optimal solutions by directing the search process toward promising areas of the search space.

In line with this, \textcite{BarbCiorbea2023} compared an exact method, Mixed Integer Linear Programming (MILP), with a metaheuristic approach, Ant Colony Optimization (ACO), for solving the TSP with time window constraints. The results of the exact approach shown is advantageous because it guarantees the best solution; however, its runtime becomes very high as the number of cities increases. For larger instances, the heuristic approach produced solutions much more quickly, although without the guarantee of global optimality. Despite this, both algorithms were able to achieve the best result in the tested scenario. Moreover, the metaheuristic demonstrated an additional advantage by aiming to minimize waiting duration at nodes when waiting could not be avoided. These findings emphasize that while exact methods provide precision, metaheuristic methods are better suited for larger and more time-sensitive problems \parencite{BarbCiorbea2023}. This demonstrates that while exact methods are strictly optimal with respect to accuracy, their exponentially growing time complexity renders them impractical to apply to large-scale or time-constrained situations.

The goal of most metaheuristic algorithms is to find an optimal path to an instance of the TSP. The increased robustness was also found in the study of \textcite{Hossain2024} where they compared new optimization algorithms against old optimization algorithms for solving the TSP and concluded that the new algorithms demonstrated greater effectiveness than the old algorithms in medium-scale instances. The new algorithms used in the study which are the Artificial Bee Colony (ABC), Grey Wolf Optimization (GWO), and the Salp Swarm Algorithm (SSA) performed better routes on average compared to the classic algorithms such as the Genetic Algorithm (GA), Ant Colony Optimization (ACO), and Simulated Annealing (SA). This means that on average, the three newer algorithms produced more consistent results than the classic three. However, they also stated that although newer algorithms produced better outputs, traditional algorithms such as GA and ACO remained to demonstrate strengths for specific instances. In terms of platform efficiency, despite the performance of newer and complex algorithms outperforming traditional methods in accuracy and even robustness, it is worth noting that their higher computational costs, like time taken, does pose specific problems within a mobile device’s resource-constrained environment.

In contrast to this, \textcite{Wadi2025} highlights the capability of ACO and Particle Swarm Optimization (PSO) in their comparative study to provide fast execution speeds albeit with minimal reduction in solution quality, with PSO being sensitive in its parameters. They also examined another swarm-based approach in the form of Elephant Herding Optimization (EHO), which consistently outperformed both ACO and PSO by achieving lower optimal costs while maintaining short execution times. Overall, their study reveals the practicality of using swarm-based approaches for problems of large-scale instance as compared to exact methods. And in terms of route planning, algorithms that focus more on faster execution time despite lesser quality solutions prove to be more suitable for environments facing hardware-related limitations such as the processing power of mobile devices.

In the same perspective, many nature-inspired optimization techniques are still being developed in the present times. Some of these nature-inspired techniques are also related to route planning such as the Hovering Scouts and Foraging Flocks Pied Kingfisher Optimizer (HSFFPKO), an algorithm combining the known capability of Pied Kingfisher Optimizer (PKO) with two more mechanisms, Hovering Scouts and Foraging Flocks, to balance its power by introducing another way of exploration and exploitation when searching for a solution. In this case, the Hovering Scout component provides PKO with an increased early-stage exploration through probing of new regions of solution space, then the Foraging Flocks mechanism serves as the improved exploitation component by guiding candidate solutions towards better areas and thereby improving convergence accuracy. HSFFPKO showed strong performance in optimization on the CEC 2017 benchmark suite and on various engineering design problems by showing faster convergence and high solution quality compared with several swarm-based metaheuristics \parencite{DelaCruz2025HSFFPKO}.

Each discussed metaheuristic algorithm demonstrates their strengths and caveats. For instance, in the study of \textcite{Hossain2024}, new algorithms (ABC, SSA, GWO) generally have better convergence and scalability than older algorithms (GA, ACO, SA). In terms of speed and accuracy, SA and PSO show fast runtime but sacrifice accuracy. However, the algorithms in the study of \textcite{Hossain2024} and \textcite{Wadi2025} were tested under different tests. Hence, to find the best algorithm among them, each must be evaluated under the same test. Furthermore, the review also shows that the exact algorithm is impractical due to the time it takes to find the exact solution, thus excluding it in the algorithm to be evaluated. The best performing algorithm in terms of speed and accuracy was the algorithm implemented in the study’s developed mobile application.


\subsection{Itinerary Recommendation and Generation}

An itinerary is defined as a structured sequence of selected places, also called Point of Interest (POI), from a set of destinations in an area. Itineraries often include activities that are offered by the POI. Every itinerary is planned accordingly under many constraints, especially in scheduling the order of visitation. These constraints include but are not limited by the route distance, time-allocation, and popularity. By planning an itinerary, tourists can utilize it to serve as a practical roadmap for their travels. These itineraries are not only useful to tourists, but are also an important application of what is called optimization algorithms and recommendation systems. Researchers often describe it as a tourist trip design problem (TTDP), a combinatorial planning task that finds the optimal visiting order of POIs under many constraints \parencite{Ruiz-Meza2022systematic}.

A recommendation system is often included in itinerary generators as explored in the study of \textcite{Otaki2025} where they create a Travel recommender systems (TRSs) that recommend the most relevant itineraries for the users. One of the popular recommendation methods is the Content-based Filtering (CBF). It uses a feature list of items and compares it with items preferred by a specific user previously. The items that match in similarity are recommended to the user. CBF works by storing user profiles based on item features which are most commonly preferred by the user. These features are used to map the similarity of one item with another by similarity equation. Then, it compares each item’s features with the user profile and recommends it based on the degree of similarity \parencite{Raghuwanshi2019}.  The other popular recommendation method is through Collaborative Filtering (CF). This method is used to predict a user’s preferences or opinions by using the collective information of other users of the system. It basically analyzes the similarity between user’s preferences and provides a recommendation that way \parencite{Li2023,Widayanti2023}. However, both CBF and CF suffer from a cold start problem where it requires a large enough data of users.

Another approach for itinerary recommendation is based on the Orienteering Problem (OP). According to \textcite{Lim2019} Tour recommendation has its roots in the OP and similar variants where a key feature is that they do not incorporate any personalization for individual users. However, a study by \textcite{Yochum2020} proposed an Adaptive Genetic Algorithm with dynamic crossover and mutation probabilities (AGAM) to handle personalized multi objective itinerary planning. Their approach integrates mandatory POIs, popularity ratings, visit duration, travel time, and costs into a weighted fitness function. By using data from platforms like TripAdvisor, GoogleMaps, and Flickr, AGAM adapts during evolution to avoid stagnation and produce diverse solutions. They tested the model with a dataset for six cities namely Budapest, Edinburgh, Toronto, Glasgow, Perth, and Osaka. AGAM was compared against two baseline heuristics, MaxN a greedy approach that prioritizes the maximum number of POIs, and MaxP also a greedy approach that prioritizes POIs with the highest rating. The results show that while MaxN and MaxP performed better at including mandatory POIs, AGAM generated richer itineraries with higher-rated POIs. AGAM shows more effective use of time budgets, and improved user enjoyment, but at the expense of higher cost and fewer mandatory POIs. This study shows that an algorithmic approach can provide personalized itinerary for tourists.

The algorithmic approach of \textcite{Yochum2020} highlights adaptability through genetic algorithms and diverse real world data sources to produce personalized itineraries and AGAM prioritizes richer and higher rated itineraries even at the cost of efficiency. Additionally, this approach solves the shortcomings of CBF and CF such as: CBF doesn’t have much control over route-level constraints like travel time, distance and budget. Due to the shortcomings of CF and CBF and since the proposed system is a standalone mobile application with no external server, an algorithmic approach is the desired approach for TRS. Hence, the algorithmic approach of \textcite{Yochum2020} is the framework that is implemented on the developed system because it combines recommendations and itinerary generation through the AGAM framework.

\subsection{Mobile Itinerary Generator}

Itinerary generators have been adopted in mobile devices, and various studies have consistently highlighted the role of smartphones in automating and personalizing trip planning. A study by \textcite{Estilo2023} developed a mobile application named Destamp, it is a comprehensive planning tool that unifies itinerary generation, budget tracking, and recommendations for Iloilo city. Destamp uses Genetic Algorithm combined with CBF for recommendations. \textcite{Loh2022} also developed an itinerary planner with functionalities like trip generation, albeit leaning more to being a social platform taking perspective in tourism. Leong’s application consists of itinerary-related functionalities such as itinerary management, route building function, and an optionally available manual editing of the said itinerary. Another application developed by \textcite{Rahman2025} called iGuide takes the angle of smart-tourism solutions for heritage cities by integrating both itinerary management with location-based services (LBS), in which they embed a real-time interactive map for promoting local attractions.

Personalization is another thing that makes an itinerary relevant to travellers. \textcite{Yulfihani2024} developed a recommendation system for the tourism sector of Batang Regency in Indonesia. Their developed system is based on a system that compares categorical data of destinations then pairs it with user preferences, which makes the recommended itinerary better fit to the user’s interest. For \textcite{Estilo2023} their user preference takes into account financial constraints in addition to the choice of attractions within Iloilo City. On the other hand \textcite{Rahman2025} and iGuide Ipoh’s personalization aspect heavily focused on the contextual and spatial personalization through LBS. iGuide Ipoh helps users discover attractions in a proximity-based approach then suggests it in real-time to the user to manage their trips while in the current destination. This approach however, relies less on sophisticated preference modeling, and more on location and maps with curated city contents. The personalization feature of the application of \textcite{Loh2022} offers it through categorized attractions, and user-defined trip parameters such as dates, starting and ending points, and duration, which then generates a route which can still be rearranged by users. This approach balances automation and user-control over the itinerary.

These studies complement each other, for instance the study of \textcite{Estilo2023} lacks itinerary management, which \textcite{Loh2022} and \textcite{Rahman2025} included in their study. Additionally, \textcite{Yulfihani2024}, although providing a great tourism recommendation system, didn’t include interactive map capabilities just like the other mentioned studies. Lakad sought to implement these features to fill the shortcomings with each mentioned study. By unifying personalized recommendations, route optimization, itinerary management, and map-based navigation; Lakad positions itself as an improved and more holistic mobile solution for travel planning.


\section{Review of Related Studies}

In this section, relevant studies regarding the developed application are presented. Covering topics related to tourism in Bulacan, techniques for optimization of itineraries and itinerary generation, along with several related approaches in personalized travel planning. The review of these studies established the foundation of the system, and the selection for the most suitable methods for mobile itinerary generation is justified.

\subsection{Trippit: An Optimal Itinerary Generator}

Trippit by \textcite{Nguyen2019} addresses a fundamental problem in itinerary planning which is the optimization of the itinerary itself which is a fundamental limitation of Google Maps. Although Google Maps supports multiple stops, it leaves the sequencing to the user and often resulting in inefficient trip planning. Manual planning takes time and significant effort to optimize in terms of total distance traveled and how long each location to visit for. The study also stressed the fact that most travelers would often switch between Google Maps and Yelp for planning itinerary. The system eliminates the difficulty in travel planning which makes vacation trips more time efficient and enjoyable to travelers.

The system was developed using the Waterfall method where each design phase is completed first before beginning the next. An android application was developed using React Native library and called network APIs such as the Google Places API and Foursquare Places API to calculate the optimized itinerary provided by the user. Users can input point of interests (POI) in which the application then retrieves the information from the respective APIs and calculates and outputs the optimized tour based on a greedy algorithm. Their system also calculates the best time to visit each input.

Trippit underwent several testing phases: primary testing, where they compared the results of the system to a known small list of itineraries; unit and integration testing, where they tested the logic and the correctness of the user interface; and alpha testing, with 10 users and was validated using a user satisfaction survey, in which most of the alpha users were satisfied with the resulting itinerary.  However, through the selection of the algorithm  and device came its limitation, which is that the application can only process up to 12 points of interest (POI). Its evaluation is also a concern as it didn’t use any standard evaluation methods like ISO/IEC 25010:2023 standard or the Technology Acceptance Model (TAM). Although a small prototype, Trippit can be considered a good enough system as it was able to produce good itineraries with a simple smartphone application without the use of advanced tools like Google OR.

However, the system focuses solely on itinerary optimization and does not include any recommendation features. Users must already know which POIs they want to visit and input them into the system. There is no mechanism to suggest POIs based on user preferences or popular attractions. Additionally, the greedy algorithm used may not always yield the globally optimal itinerary, especially as the number of POIs increases. The limitation of processing only up to 12 POIs also restricts its applicability for longer trips with more destinations. Future improvements could include integrating recommendation algorithms to suggest POIs and exploring more advanced optimization techniques to handle larger sets of locations.

\subsection{Optimization of Tourism Destination Recommendations in Batang Regency Using Content-based Filtering}

\textcite{Yulfihani2024} created a tourism recommendation system for Batang Regency in Indonesia. They applied Content-Based Filtering (CBF) for their recommendation system, using tourist preferences as their data source. Tourists interact with the system by adding destinations to a wishlist. The system then uses the wishlist to learn the tourist’s preferences. The recommendation factors category similarity where it compares the categories of wishlist items like nature, culture, culinary, and leisure; and location similarity by calculating the distance between destinations using the Haversine formula, which finds the shortest distance between two points on Earth using latitude and longitude.  The total similarity then calculated with 70 percent category similarity and 30 percent location similarity. The system ranks all available attractions by their total similarity score and the top 10 destinations are recommended to the tourist. The system was tested in different scenarios such as single item in wishlist, multiple items from same category with same/different locations, and multiple items from mixed categories. They use precision, recall, and F1 score as a metric for the system.

The recommendation system was designed with a three-tier architecture that includes Backend, Web, Admin Panel, and Mobile Application. The system’s backend was developed using Laravel, under the PHP Framework, and MySQL. The backend part of the system handles all the operations including but not limited to operations relating to user authentication and processing recommendations of the CBF Algorithm. The backend is also built to manage the communication between the mobile application, the database, and the admin panel. While the backend manages backbone tasks, the Web Admin Panel is built for administrators. The web admin is still developed using Laravel and enables admins to manage tourist destinations and its categories. Changes made by administrators in the web admin panel will update the database through the backend of the system, instantly reflecting the updated information to the mobile application. The mobile application side of the system was developed using Android Studio and the Kotlin programming language. This mobile application side of the system provides end-users with the login and registration pages, along with a wishlist management page where the personalized recommendation is based on using CBF. It also includes additional features such as attraction details panel, and filtering by categories and proximity.

Overall, the system shows highly accurate and relevant recommendations for tourists, especially in simple preference scenarios. The F1 result of 0.965 also shows that the CBF approach works well when aligning recommendations with tourists’ interests. Its performance however, dropped slightly when complex scenarios arise, issues such as mixed categories and location within the wishlist items. \textcite{Yulfihani2024} recommended to future researchers to expand the dataset and use hybrid models like combining CBF with Collaborative Filtering. They also suggested that adding user feedback mechanisms may be able to improve the relevance of recommendations.

This case of Batang Regency focused heavily on the recommendation part of itinerary generation. Despite being proven to be a robust and good system in POI recommendation, it lacks the capability of providing tourists with an optimized routing for a much more efficient itinerary optimization. Meaning that while the system can recommend destinations based on user preferences, it does not optimize the sequence of visits or consider travel constraints like time and distance. Additionally, they didn’t include an evaluation for the system like ISO/IEC 25010:2023 standard or the Technology Acceptance Model (TAM), only the evaluation for
content-based filtering was provided.

\subsection{Visit Planner: A Personalized Mobile Trip Design Application based on a Hybrid Recommendation Model}

Visit Planner (ViP) a mobile application prototype developed by \textcite{papadakis2024}. It offers personalized recommendations for an itinerary based on user preference which are either explicitly collected by the application or assessed through the user’s behavior within it. ViP utilizes an Expectation Maximization method to offer the user an itinerary that tailors to their satisfactory needs while taking into consideration time and spatial constraints that concern both the user and the destination. The application currently focuses on the city of Agios Nikolaos in Crete. The system requires a user to create a secured profile in a registration process, wherein demographic information and essential data for the algorithm is collected. After this registration process, the user is also asked to specify their preferences in three different ways of their choice, one is by rating categories of POI, another choice is by stating if they like or dislike the category, and another one is by selecting their most preferred category. These categories are carefully specified by analyzing responses collected from 150 visiting tourists, in addition to local and expert knowledge in Agios Nikolaos.

The architecture of the system consists of main components like the front-end User Interface, back-end database, the middleware for processing the information, the recommendation components, and the itinerary creation components. The front-end of the application is an android app available in Google Play. The back-end is composed of MariaDB databases for the purpose of storing necessary information of users and POIs for smooth operations and functionalities. The middleware, built upon the basis of Spring Boot Framework aims to process the controlling flow between the system components and perform the functionalities. It works by receiving queries from the front-end, and then sending those to the recommendation algorithms. The recommendation and itinerary creation components comprises four several recommendation algorithms to tender the needs of the user, where each user is assigned a different algorithm for their cases in performance evaluation wise. ViP integrates multiple, novel recommendation algorithms which can adapt to different user profiles.

The algorithms include model-based collaborative filtering that uses synthetic coordinate based system for recommendations (SCoR), hierarchical content-based similarity measures, a hybrid content-based recommendation using Weighted Extended Jaccard approach combined with the second one, and Bayesian recommendation algorithm. Any of these algorithms is applied to generate personalized POI suggestions. And once candidate POIs are retrieved, an expectation-maximization-based itinerary creation component sequences them into a feasible route by maximizing user satisfaction function constrained under temporal and spatial factors.


\subsection{I-AIR: Intention-aware travel itinerary recommendation via multi-signal fusion and spatiotemporal constraints}

\textcite{Cui2025} developed I-AIR (Intention-Aware Itinerary Recommendation), a system based on deep learning that aimed to generate personalized and practical traveling itineraries. I-AIR’s primary components consist of what they called Itinerary Construction Policy and Intention-Aware POI Scoring Model. Itinerary Construction Policy organizes selected POIs into a logical and sequential trip. The Intention-Aware POI Scoring Model then assesses the relevance of possible POIs in the context of user preference. The user preference model of I-AIR is a combination of applied transformer-based sequential encoder to identify time patterns from check-in history, and an explicit feedback encoder that synthesizes preferences from ratings and likes, this is in addition to utilizing a Graph Convolutional Network (GCN) for representing POI relationships from patterns like co-visitation. I-AIR integrates all of these components using a multi-signal fusion layer for the generation of customized POI relevance scores. I-AIR also features a routing optimization that uses greedy algorithms to pick top-scoring POIs viable enough at each step, while also removing the possible instances of the itinerary violating user-defined constraints. Its model is trained using combinations of predicting next-POI along with an explicit feedback reconstruction from datasets that contain user trajectories, ratings, and dwell time.

The system of I-AIR was tested on eight real world datasets. These datasets include four theme parks and four city-based tourism. It includes detailed information such as user check-in records, attraction waiting time, and tourist reviews that provided the system with the implicit and explicit feedback. The evaluation was measured under precision, recall, and F1 scoring. The comparative baseline systems then included Greedy-Popular (G-Pop), Greedy-Near (GNear), PersQ, EffiTourRec, DCC-PersIRE, BERT-Trip, and DLIR, wherein I-AIR showed that it outperformed the other systems significantly across the given datasets, yielding higher precision in both theme parks and city-based tourism scenarios.

\subsection{Destamp: A Mobile Application for Generating Budget-Based Itineraries to Enhance Travel Planning}

\textcite{Estilo2023} have proposed a very similar system to the study’s proposed application where itinerary generation and itinerary optimization are the main core features. The system of \textcite{Estilo2023} is a mobile application that uses GA to generate an itinerary for the user. However, unlike Lakad, the GA used in the study is univariate and does not consider other factors that may affect the user’s experience on the generated itineraries such as the user’s categorical preferences and POI ratings. The main objective function implied in the study is the minimizing of the itinerary travel budget depending on the maximum budget provided by the user and fitting the opening hours of each POI to the specified starting time of the user.

Destamp was developed using the AGILE methodology and the React Native framework for developing cross-platform mobile applications. Mapbox service was also used to provide mapping functions such as routing and interactive map features. While the system is split into two sections for the architecture: frontend and backend. The frontend handles all the graphical user interface and interactions which sends requests to the backend for the processing and logic such as itinerary generation.

The application has many similar features to Lakad such as itinerary generation and live map interaction. Users can set multi-day itineraries as well as allocate budgets for the itinerary to be generated. Each POI in the app contains the map location, budget cost, opening hours, and contact information. Business owners can also use the app to submit their own shops and businesses to show up as POI in Destamp, making the whole POI list accessible and dynamic.

While very similar to Lakad, Destamp is only bounded to the places in Iloilo City. It also implements a freemium model to its features where most features such as maximum itineraries are locked behind a subscription. It also uses GA for its itinerary generation but is lacking the option for the itinerary optimization through TSP model as their optimization refers to pair-wise optimization through Dijkstra algorithm instead. There is also a clear gap in itinerary management where itineraries cannot be edited and must be regenerated for any changes. Furthermore, the evaluation instrument of their study is not clearly defined, and the evaluators they considered consisted of ten respondents, in which they concluded that GA actually does improve travel recommendations.

Destamp limits their users from fully utilizing their application. This is different from Lakad which allows users full access to all the available features mentioned in Destamp such as itinerary generation and optimization. The itinerary generation of Lakad is based on the AGAM framework of \textcite{Yochum2020}, and applying TSP algorithms for route optimization. Finally, users can also fully manage their own itineraries allowing them to edit the POIs included in existing itineraries.

\section{Synthesis of the Review}

The reviewed works establish valid grounds for the developed mobile application’s technical framework. Concepts such as Graph Theory, the Traveling Salesman Problem, and Orienteering Problem provides the mathematical foundation for itinerary generation and route optimization. While exact algorithms guarantee the best solution, they are often impractical for mobile devices due to high computation times. Consequently, metaheuristic algorithms have been shown to outperform exact algorithms as explored in the study of \textcite{Hossain2024} and \textcite{Wadi2025}, offering a viable solution for mobile constraints. Regarding personalization, much of the literature also suggests that CBF suffers from the "cold-start" problem despite its effectiveness. In addition to this, since CBF is naturally a deterministic approach, the recommendations it produces are the same iteration most of the time. Algorithmic approaches on the other hand, offer greater flexibility by allowing the fitness functions to take into account multiple objectives such as user-defined parameters like preferences, algorithms such as AGAM by \textcite{Yochum2020} in particular. In this case, AGAM is a more suitable method for the application of an itinerary recommendation system, as opposed to content selection methods.

The feasibility of smartphone-based route generation application is proved to be effective when  implemented as seen with mobile applications such as Trippit, Destamp, and I-AIR. There are, however, some significant fragmentations in the current landscape since several features of the mentioned mobile applications are often exclusive to each other. Trippit prioritized a simple routing technique without further personalization, Destamp focused solely on budget, and others lacked comprehensive itinerary management features. In addition to this, many of the existing studies are being limited in terms of their reliability and usability claims as they fail to apply quality evaluation in their frameworks, evaluation metrics such as ISO/IEC 25010:2023 and Technology Acceptance Model (TAM). LAKAD fills these limitations by providing an itinerary generator verified and specifically tailored for Bulacan tourism.

All of the literature, taken together, highlights the clear demand for an itinerary generator and optimization system to exist in Bulacan. And that reliance on social media leads to misinformation, despite the underlying digital strategies serving as a crucial support for promotion \parencite{Canet2025S, DelaCruz2022}. So while Bulacan's government initiated programs such as the Pamana Pass, the province currently has no means to automate and optimize personalized itinerary planning for its tourism sector. There also exists Philippine-based studies, particularly from \textcite{Estilo2023}, but they are bound and limited to certain cities, which in Estilo's case is Iloilo City, and often do not address the problem in a provincial-level scope as required by Bulacan. All of this led LAKAD to address the underlying gaps by integrating data, curated locally and verified by the provincial government, with advanced optimization algorithms to provide tourists a personalized experience in Bulacan.

\section{Conceptual Framework}

This section explains the key inputs, processes, and outcomes in development of the study using an Input-Process-Output (IPO) model as shown in Figure \ref{fig:conceptualframework}.
\begin{figure}[H]
	\centering
	\caption{Input-Process-Output (IPO) Model}
	\includegraphics[width=0.7\textwidth, trim={1cm 2cm 1cm 2cm}, clip]{IPO.png}
	\label{fig:conceptualframework}
\end{figure}

\subsection{Input}
The input contains the review of the methods used by related studies and systems, which was used as foundation in the study. These methods consist of different algorithms used for solving the TSP and selecting which algorithm is the most applicable in the mobile platform as well as which recommendation model is a fit for itinerary recommendation. The second input is the gathering of tourist spot locations or point of interests (POI). These POIs are locations in Bulacan which can be considered as tourist attractions. This was gathered from and verified by the Provincial History, Arts, Culture, and Tourism Office (PHACTO) to assess the correctness of the POIs included in the system. Development tools refer to the hardware and software tools that was used to develop the application. These include operating systems such as Windows or Linux, integrated development environments (IDE) such as Visual Studio or Android Studio, and design pieces of software like Figma or Adobe Illustrator.

\subsection{Process}
The AGILE methodology was used in the software development of Lakad, specifically the Kanban method. Kanban method is a visual process management system that can manage knowledge and work by considering the Just In Time (JIT) delivery approach \textcite{Alaidaros2021}. The first phase of the development was planning and analyzing the different requirements of the system such as flow charts, entity relationship diagrams, and use cases. Afterward, mockup user interface was made to fully outline the look and feel of the application. Building phase was the actual development of the core features of the application and testing and debugging refers to the testing of the correctness of the core features’ implementation. The power of AGILE methodology was used to iteratively cycle between building phase and testing to evaluation by users and professionals to quickly enhance features and get immediate feedback. The system went through rigorous evaluation by professionals through ISO 25010:2023 and through TAM by its intended users.

\subsection{Output}
The expected output of this study is a fully functional mobile itinerary generator and optimization system in the scope of Bulacan. Lakad application was able to optimize the route to be taken for the user’s itinerary as well as recommend personalized itineraries that may interest the user.

\section{Definition of Terms}
This section provides the definition of the terms used in the study, along with the operational aspects of the proposed system, to ensure clarity and a proper understanding of the concepts within the context of the research.

\begin{description}
	\item[Cold Start.]
	A condition in which the system lacks sufficient initial user preference data,
	limiting its ability to generate accurate personalized itinerary recommendations.
	
	\item[Distance Matrix.]
	This refers to the structured representation of pairwise travel distances or travel times between point of interests, and is used as input for route optimization algorithms in the system.
	
	\item[Pasalubong Centers.]
	This is the term referring to local establishments in Bulacan that offer regional products and souvenirs, it is integrated into the system as optional stops recommended based on proximity during route traversal of itineraries.
	
	\item[Personalized Itinerary Recommendation.]
	Refers to the system-generated travel plan that selects and orders destinations in accordance to user-defined preferences, constraints, and optimization results.
	
	\item[Provincial History, Arts, Culture, and Tourism Office (PHACTO).]The official government agency that oversees the tourism sector of Bulacan, they serve as the primary source of verified points of interest that are included in the system.
	
	\item[Point of Interest (POI).]
	These refer to the verified tourist destinations within Bulacan, and are provided by and contains information validated by PHACTO, these serve as the candidate nodes for tourist spot suggestions.
	
	\item[Tourist Spot Searching.]
	Refers to the feature of the system that allows users to query and retrieve available points of interest included in the database, which displays the information about the tourist spot or POI.
	
	\item[Mapbox.]
	Refers to the third-party mapping service that is integrated into the system for providing the interactive map, geolocation visualization, and navigation support for the generated itineraries.
	
	\item[OpenRouteService (ORS).]
	Refers to the routing service utilized by the system to compute travel
	distances, travel duration, and the optimized paths between selected points of interest.
	
	\item[Application Programming Interface (API).]
	Are a set of protocols that enables the developed system to communicate with the backend, like Mapbox and OpenRouteService, to retrieve necessary functionalities like mapping, routing, and obtaining location data.
	
\end{description}
